\section{Exploratory Analysis}
To understand the underlying structure and distribution of our dataset, we conducted a detailed exploratory data analysis. The analysis includes summary statistics, visualizations (scatterplots, heatmap), and correlation analysis.
\subsection{Descriptive Statistics}
We have computed basic descriptive statistics (mean, median, range, SD) for each of the variables like Rainfall (mm) , Avg Temp (\degree C) , Fertilizer (kg/ha), Pest Incidents , Technology Index , 	Crop Yield (ton/ha). Now let us have a look into it.


\begin{figure}[h!]
    \centering
    \includegraphics[width=1\linewidth]{Screenshot 2025-04-15 134414.png}
    \caption{Basic statistics summary}
    \label{fig:enter-label}
\end{figure}

Looking into the above table the following basic things we can say about the data set : 
\begin{itemize}
   

 \item Moderately high variability in all of the variables implies that none of the variables are remotely stable with time.

 \item Pest incident is the only positively skewed factor(apart from temperature) suggesting the variability is caused mainly by few years when the pest incidents are at their peak.


 \item High range in Technology Index shows the improvement of its implementation in agriculture with time.


 \item Rainfall varies significantly with a high range and being an important factor in crop yield is worth studying.


 \item Temperature on the other hand has a lower variance and range compared to Rainfall. So it can definitely be considered a more stable natural factor. 
 
\end{itemize}
\subsection{Correlation with Crop Yield}

First have a look into the result for Pearson correlation coefficient between crop yield and all other variable one by one. 
\begin{figure}[h!]
    \centering
    \includegraphics[width= 1.2\linewidth]{Screenshot 2025-04-15 142653.png}
    \caption{Correlation}
    \label{fig:enter-label}
\end{figure}

\textbf{\underline{INSIGHTS}}:- 

\vspace{1 mm}
\begin{itemize}
\item Very high negative correlation between Temp and Yield implies temperature and global warming as a important factor affecting agriculture

\item Moderately high and positive linear link between Crop yield and Fertlizer used as well as Technology suggest that these two factors require proper attention for maximising agricultural progress.

\item A high positive correlation with Rainfall tells us importance of rainfall, the main source of irrigation in agriculture
\end{itemize}

\subsubsection{Correlation Matrix}


\begin{figure}[h!]
    \centering
    \includegraphics[width=1\linewidth]{Screenshot 2025-04-15 143434.png}
    \caption{Correlation Matrix}
    \label{fig:enter-label}
\end{figure}

\subsubsection{Heat Map}
\begin{figure}[h!]
    \centering
    \includegraphics[width=1\linewidth]{Screenshot 2025-04-15 145845.png}
    \caption{Correlation Heat Map}
    \label{fig:enter-label}
\end{figure}

\textbf{\underline{INSIGHTS}}:- 

\vspace{1 mm}
\begin{itemize}
    
\item Rainfall seems to have a strong positive correlation with both Fertilizer used and Technology Index in the agricultural model. The former might be because more rainfall implies more waste of fertlizers by leaching or run-off.

\item Fertilizer and Technology Index both has a high negative correlation with Pest Incidents which is pretty self-explanatory.

\item Technology Index and Fertilizer has a very strong positive correlation which suggests that farmers tend to invest in Fertilizers and Technology side by side and improving them almost equally.

\item Rainfall and Temperature has a moderately high negative correlation which is expected.
\end{itemize}
\subsection{Scatter Plot}

Now we are looking into the following scatter plot : -

\vspace{1 mm}
\hspace{-5 mm}i) Scatter plot for Yield Vs Rainfall

\vspace{1 mm}
\hspace{-5 mm}ii) Scatter plot for Yield Vs Avg Temp

\vspace{1 mm}
\hspace{-5 mm}iii) Scatter plot of Yield vs Fertilizer

\vspace{1 mm}
\hspace{-5 mm}iv) Scatter plot for Yield Vs Pest incidents

\vspace{1 mm}
\hspace{-5 mm}v) Scatter plot for Yield Vs Technology index

\vspace{1 mm}
\hspace{-5 mm}We then discuss some basic insights from the scatter plots.

\vspace{10 mm}
\begin{figure}[h!]
    \centering
    \includegraphics[width=0.75\linewidth]{Screenshot 2025-04-15 151149.png}
    \caption{Scatter plot of Yield vs Rainfall}
    \label{fig:enter-label}
\end{figure}
\textbf{\underline{INSIGHTS}}:- 

 \begin{itemize}
 
 \item As rainfall increases, the crop yield also increases, which is shown by the upward trend and the fitted blue regression line.

\item This suggests that rainfall plays a significant role in boosting agricultural productivity — likely because water is a key input in crop growth.
\end{itemize}

\newpage
\begin{figure}[h!]
    \centering
    \includegraphics[width=0.75\linewidth]{Screenshot 2025-04-15 152707.png}
    \caption{Scatter plot of Yield vs Avg Temp}
    \label{fig:enter-label}
\end{figure}

\textbf{\underline{INSIGHTS}}:- 

\begin{itemize}
\item The blue regression line slopes downward, indicating that higher temperatures are associated with lower yields.

\item This makes sense agriculturally:

i) High temperatures can cause heat stress in crops.

ii) It increases evapotranspiration, drying out soil moisture faster.
\end{itemize}

\begin{figure}[h!]
    \centering
    \includegraphics[width=0.75\linewidth]{Screenshot 2025-04-15 153249.png}
    \caption{Scatter plot of Yield vs Fertilizer}
    \label{fig:enter-label}
\end{figure}

\textbf{\underline{INSIGHTS}}:- 

\begin{itemize}
    
\item The blue regression line shows an upward trend, indicating that increasing fertilizer tends to increase yield.

\item However, there's more scatter around the line compared to the rainfall plot, which means:

i) The effect of fertilizer on yield is less consistent.

ii) Overuse of fertilizer could lead to negative effects, possibly explaining some of the outliers.
\end{itemize}

\begin{figure}[h!]
    \centering
    \includegraphics[width=0.75\linewidth]{Screenshot 2025-04-15 153614.png}
    \caption{Scatter plot of Yield vs Pest incidents}
    \label{fig:enter-label}
\end{figure}
\textbf{\underline{INSIGHTS}}:- 

\vspace{1 mm}

\begin{itemize}
\item The blue regression line slopes downward, confirming the overall negative trend.

\item Pest control alone isn't sufficient to ensure high productivity.
\end{itemize}

\begin{figure}[h!]
    \centering
    \includegraphics[width=0.75\linewidth]{Screenshot 2025-04-15 154617.png}
    \caption{Scatter plot of Yield vs Technology Index}
    \label{fig:enter-label}
\end{figure}

\textbf{\underline{INSIGHTS}}:- 
\vspace{1 mm}

\begin{itemize}
\item The blue regression line slopes upward, showing a strong positive linear trend.

\item Investment in agricultural technology (like irrigation systems) pays off in terms of higher yield.
\end{itemize}

\subsection{Partial Correlation}
\begin{figure}[h!]
    \centering
    \includegraphics[width=0.75\linewidth]{Screenshot 2025-04-15 144100.png}
    \caption{Partial correlation with crop yield}
    \label{fig:enter-label}
\end{figure}

\textbf{\underline{INSIGHTS}}:- 

\vspace{1 mm}
\begin{itemize}
    \item Controlling temperature, partial correlation between yield and rainfall is 0.390 which is significantly less than their linear correlation. That can be explained as temperature and rainfall has a moderately negative correlation between them.

\item Now keeping rainfall constant and doing it with temp, the value is -0.875 which is a slight decrease in magnitude from their individual correlation. This leads to the conclusion that temperature is a more important factor than rainfall if we want to explain the linear variability of Crop yield subject to linear variation of these variables.

\item Technology Index and Fertilizers being highly correlated, we did the same partial correlation calculation with them too.

\item Technology and Crop Yield controlling Fertilizer shows almost no correlation ie linear independence. So it might be that they are related with some non-linear function. On the other hand the correlation strength between Fertilizer and Yield shows a moderate decline which tells Technology has some effect in the initial correlation between them.

\end{itemize}
\subsection{ Multiple Linear Regression with Crop Yield as dependent variable and others as independent variable}

\begin{figure}[h!]
    \centering
    \includegraphics[width=0.65\linewidth]{Screenshot 2025-04-15 180013.png}
    \caption{Residuals(deviation from fit)}
    \label{fig:enter-label}
\end{figure}
\begin{figure}[h!]
    \centering
    \includegraphics[width=0.75\linewidth]{Screenshot 2025-04-15 173216.png}
    \caption{Multiple Linear Regression}
    \label{fig:enter-label}
\end{figure}

\begin{itemize}
  \item \textbf{Multiple R-squared:} 0.9339
  \item \textbf{Adjusted R-squared:} 0.9165
\end{itemize} 



\subsection{Time series graph}
We have plotted three time series graph for total crop yield , amount of rainfall , avg temperature over the years from 2000 to 2024. Let's have a look in it and discuss what can we say from these graphs.


\begin{figure}[h!]
    \centering
    \includegraphics[width=0.67\linewidth]{Screenshot 2025-04-11 143416.png}
    \caption{Time series plot for crop yield}
    \label{fig:enter-label}
\end{figure}

\subsubsection{Some Observations from the Crop Yield time series graph}

1) \textbf{\textcolor{blue}{High Variability}:} Crop yield fluctuates significantly from year to year, ranging between approximately 3 to 10 tons per hectare, indicating sensitivity to varying conditions.

\vspace{2 mm}
\hspace{-5 mm}2) \textbf{\textcolor{blue}{No Clear Trend}: }There is no obvious long-term increase or decrease. However, the yields show short-term peaks and dips, possibly linked to environmental or agricultural factors.

\vspace{2 mm}
\hspace{-5 mm}3) \textbf{\textcolor{blue}{Low Yield Years}}: Major drops are seen around 2009, 2012, 2014, 2020, and 2023, where yields fall close to or below 4 tons/ha.

\vspace{2 mm}
\hspace{-5 mm}4) \textbf{\textcolor{blue}{High Yield Years}:} Peaks occur around 2001, 2004, 2008, 2013, and 2024, with 2024 marking one of the highest yields, nearly reaching 10 tons/ha.

\vspace{2 mm}
\hspace{-5 mm}5) \textbf{\textcolor{blue}{Potential Recovery in 2024}:} Despite previous dips, crop yield in 2024 shows a strong recovery, suggesting improved conditions or interventions.

\begin{figure}[h!]
    \centering
    \includegraphics[width=0.67\linewidth]{Time series of avg Temp.png}
    \caption{Time series plot for avg temp}
    \label{fig:enter-label}
\end{figure}

\subsubsection{Some Observations from the avg temp time series graph}

\vspace{2 mm}
1) \textbf{\textcolor{blue}{Fluctuating Pattern}:} The average temperature varies significantly year by year, ranging mostly between 21°C and 35°C, showing no consistent upward or downward trend.

\vspace{2 mm}
\hspace{-5 mm}2) \textbf{\textcolor{blue}{Highs and Lows}:}
Peaks around 2005, 2012, 2015, 2017 and 2021 (close to or slightly above 35°C).
Lows around 2004, 2008, 2013, 2019, and 2024 (dropping near or below 22°C).

\vspace{2 mm}
\hspace{-5 mm}3) \textbf{\textcolor{blue}{No Strong Warming Trend}: }Despite climate concerns, this data does not show a steady increase in temperature over time—variability dominates instead of a clear warming signal.

\vspace{2 mm}
\hspace{-5 mm}4) \textbf{\textcolor{blue}{Sharp Drop in 2024}: }The year 2024 recorded one of the lowest average temperatures in the entire dataset.

\begin{figure}[h!]
    \centering
    \includegraphics[width=0.67\linewidth]{Time series of rainfall.png}
    \caption{Time series plot for rainfall}
    \label{fig:enter-label}
\end{figure}

\subsubsection{Some Observations from the rainfall time series graph}
1) \textbf{\textcolor{blue}{Fluctuating Trend}:} Rainfall shows significant year-to-year variation, indicating a highly fluctuating pattern without a clear long-term increasing or decreasing trend.

\vspace{2 mm}
\hspace{-5 mm}2) \textbf{\textcolor{blue}{No Strong Trend}: }There is no consistent upward or downward trajectory, which means rainfall has not significantly increased or decreased over the 25-year period.

\vspace{2 mm}
\hspace{-5 mm}3) \textbf{\textcolor{blue}{Peaks and Lows}:}

    Rainfall peaked around 2001, 2008, 2011,2013, 2016, and 2024  (values around 1400–1500 mm).

    The lowest points were around 2004, 2009, 2014, and 2020 (values around 700–800 mm).
